% main.tex
\documentclass[11pt]{article}

\usepackage[margin=1in]{geometry}
\usepackage{amsmath,amssymb,amsfonts}
\usepackage{graphicx}
\usepackage{booktabs}
\usepackage{microtype}
\usepackage{hyperref}
\usepackage{url}
\usepackage{enumitem}
\usepackage{algorithm}
\usepackage{algpseudocode}
\usepackage[round]{natbib}

\title{Quantum-Assisted Hypergraph Urban Sensing (QAHUS):\\
Edge-Efficient Vision, Spatiotemporal Hypergraph Forecasting, and Constrained Decision Support}
\author{Team Archive200 \\
\small (Names omitted for blind review / competition submission)}
\date{June 2025 (revised and technicalized)}

\begin{document}
\maketitle

\begin{abstract}
We propose a computationally efficient, deployable framework that links \emph{edge vision} to
\emph{city-scale forecasting and decision support} via a \emph{dynamic hypergraph} representation of urban systems.
The core idea is to (i) extract low-latency visual signals on resource-constrained devices, (ii) fuse them with
heterogeneous civic data streams into a time-evolving hypergraph, (iii) forecast and detect anomalies using
spatiotemporal hypergraph learning, and (iv) plan interventions under explicit safety, fairness, and privacy
constraints. The quantum contribution targets two near-term bottlenecks: (a) fast similarity computation via
quantum feature maps / kernels for anomaly detection and regime classification, and (b) quantum-assisted
optimization (QUBO/QAOA/annealing) for constrained resource allocation. The overall system is designed to be
auditable, privacy-preserving (via differential privacy where appropriate), and compatible with modern
governance requirements.
\end{abstract}

\section{Project Description}

\subsection{Elevator Pitch}
We are developing a computationally efficient vision-to-decision framework that transforms how machines
perceive and respond to real-world environments. By combining high precision, seamless scalability, and
real-time adaptability, the system enables intelligent visual processing that can evolve from low-power
edge devices to advanced autonomous systems. Deployed across healthcare, robotics, environmental monitoring,
and smart-city operations, the framework bridges the gap between accuracy and responsiveness.

\subsection{High-Level Description}
Urban environments and critical domains require systems that can (i) see locally with low latency, (ii) fuse
multiple data sources reliably, and (iii) recommend actions under constraints. We unify these requirements
with a pipeline:

\begin{enumerate}[leftmargin=*]
\item \textbf{Edge vision}: on-device detection/tracking/segmentation with mobile architectures
(e.g., MobileNet-family backbones) to produce structured counts, trajectories, and events.
\item \textbf{Urban hypergraph fusion}: represent entities (intersections, facilities, sensor clusters, cohorts)
as nodes, and \emph{multi-way relations} (e.g., ``this set of sensors jointly characterizes this corridor at time $t$'')
as hyperedges.
\item \textbf{Spatiotemporal learning}: predict flows/risks and detect anomalous regimes using hypergraph neural
networks adapted to temporal forecasting.
\item \textbf{Constrained decision support}: recommend interventions using explicit constraints (safety,
service-level, equity), solved via primal--dual methods and/or safe policy learning.
\item \textbf{Quantum assistance}: accelerate key subroutines for similarity and combinatorial optimization on
near-term quantum hardware.
\end{enumerate}

\subsection{Innovation (What is New Here?)}
The novelty is \emph{not} a single algorithmic trick, but a coherent, deployable architecture that:
(i) runs end-to-end under tight compute budgets,
(ii) uses a hypergraph (not just pairwise graphs) for multi-agent, multi-sensor relations,
(iii) encodes \emph{governance constraints} as first-class mathematical objects, and
(iv) introduces quantum subroutines where they plausibly help (kernel evaluation and discrete optimization).

\section{Alignment with UN Sustainable Development Goals (SDGs)}
The system supports SDG 9 (Industry, Innovation, and Infrastructure) through efficient, low-cost deployment;
SDG 11 (Sustainable Cities and Communities) through real-time monitoring and planning; SDG 13 (Climate Action)
through rapid detection of environmental risks; and SDG 3 (Good Health and Well-being) through assistive
monitoring and diagnostics.

\section{QAHUS Technical Framework}

\subsection{Dynamic Urban Hypergraph Model}
Let the city at time $t$ be a weighted hypergraph
\[
\mathcal{H}_t = (V, \mathcal{E}_t, w_t),
\]
where $V$ indexes entities (locations, facilities, sensor regions), and each hyperedge
$e \in \mathcal{E}_t$ connects a \emph{set} of vertices that jointly participate in an event, constraint, or latent factor.
Let $x_t(v) \in \mathbb{R}^d$ be the feature vector for node $v$ at time $t$ (from edge vision, IoT, mobility, etc.).
Hyperedge features can be represented as pooled statistics
\[
x_t(e) = \mathrm{pool}\{x_t(v): v\in e\},
\]
with learned pooling (attention) or fixed pooling (mean/max) depending on deployment constraints.

Hypergraphs naturally represent multi-way dependencies that pairwise graphs often approximate poorly
(e.g., one incident involving multiple intersections and transit nodes simultaneously). We build on
hypergraph neural network ideas, using a hyperedge-induced convolution operator. \citep{feng2019hypergraph}

\subsection{Edge Vision Layer}
We assume cameras and/or low-power vision sensors provide raw frames locally.
A mobile network (e.g., MobileNetV2-style inverted residual blocks) yields on-device detections and tracks,
and outputs a compact event stream (counts, queues, trajectories, near-miss alerts). \citep{sandler2018mobilenetv2}

This output is \emph{intentionally compressed}: the system is designed to avoid shipping raw video whenever
possible, improving privacy and reducing bandwidth.

\subsection{Spatiotemporal Forecasting on Hypergraphs}
For forecasting, we adapt the spatiotemporal graph learning literature to hypergraphs.
Graph baselines include DCRNN and STGCN; QAHUS generalizes these to multi-way relations. \citep{li2018dcrnn,yu2018stgcn}

A simple (and deployable) template is:
\[
h_{t+1} = \mathrm{GRU}\Big(h_t,\ \mathrm{HyperConv}(x_t, \mathcal{H}_t)\Big),
\quad
\hat{y}_{t+\Delta} = g(h_{t+1}),
\]
where $\mathrm{HyperConv}$ is a hypergraph convolution (e.g., via incidence-based message passing), and $g$
produces forecasts (flows, occupancies, risk scores). In practice, we also use temporal convolutions or
Transformers when compute budgets allow.

\subsection{Consistency Constraints via Sheaf-Theoretic Data Fusion (Optional but Powerful)}
Many urban deployments suffer from \emph{inconsistent} multi-sensor reporting (missingness, calibration drift).
Cellular sheaves provide a principled way to encode \emph{local consistency maps} and to measure global
inconsistency via cohomological obstruction. \citep{curry2014sheaves,hansen2019spectral}

Operationally: attach measurements to cells (nodes/edges/hyperedges) and enforce linear restriction maps
that express ``what these two sensors should agree on.'' A large inconsistency score triggers recalibration,
down-weighting, or conservative decision modes.

\section{Quantum Contribution}

\subsection{Quantum Feature Maps and Kernels for Regime Classification}
We use quantum feature maps to compute similarities
\[
K(x,x') = \left|\langle \phi_\theta(x) \mid \phi_\theta(x') \rangle\right|^2,
\]
where $\lvert \phi_\theta(x)\rangle$ is a parameterized circuit embedding of classical features.
These kernels can be used for anomaly detection, regime classification, and hybrid models that combine
classical encoders with quantum similarity layers. \citep{havlicek2019supervised,schuld2019qkernel}

We emphasize that QAHUS does \emph{not} require quantum advantage to be useful: the kernel module is a swap-in
component evaluated against classical approximations.

\subsection{Quantum-Assisted Constrained Resource Allocation}
Many civic interventions (dispatching responders, rerouting flows, placing temporary controls) reduce to
combinatorial optimization with constraints. We encode a decision vector $z\in\{0,1\}^m$ and minimize an
energy
\[
\min_{z\in\{0,1\}^m}\; z^\top Q z \quad \text{s.t.}\quad Az \le b,
\]
which can be converted to a penalized QUBO and solved with quantum annealing or QAOA-style heuristics,
then verified and refined classically. The constrained formulation and dual methods follow standard convex
optimization principles. \citep{boyd2004convex,rubinstein2004crossentropy}

\section{Ethics, Safety, and Governance as Mathematical Constraints}

\subsection{Constrained Decision-Making}
We treat governance requirements as explicit constraints: safety caps, minimum service levels, and fairness
constraints over affected groups/regions. A constrained MDP viewpoint is natural:
maximize utility while keeping cumulative constraint costs below thresholds. \citep{altman1999cmdp,achiam2017cpo}

\subsection{Privacy}
We avoid raw-video centralization where possible and apply differential privacy to aggregated releases
(e.g., publicly shared counts or mobility summaries). \citep{dwork2006dp,dwork2014foundations,abadi2016dpdl}

\subsection{Fairness and Disparate Impact}
We evaluate interventions for group-level performance disparities (e.g., unequal false positive rates in
enforcement-like settings, or unequal service delays). We use standard fairness criteria such as equality
of opportunity and we explicitly audit for disparate impact risks. \citep{hardt2016equality,barocas2016bigdata}

\section{Algorithms}

\begin{algorithm}[t]
\caption{QAHUS Processing Pipeline (Deployable Template)}
\begin{algorithmic}[1]
\Require Sensors/vision frames $I_t$, civic data streams $s_t$, hypergraph schema $\Pi$
\Ensure Forecasts $\hat{y}_{t+\Delta}$, anomalies $a_t$, recommended actions $z_t$
\State $u_t \leftarrow \mathrm{EdgeVision}(I_t)$ \Comment{detections/tracks/events}
\State $x_t \leftarrow \mathrm{Fuse}(u_t, s_t)$ \Comment{compact node/hyperedge features}
\State $\mathcal{H}_t \leftarrow \mathrm{BuildHypergraph}(x_t;\Pi)$
\State $c_t \leftarrow \mathrm{ConsistencyCheck}(\mathcal{H}_t)$ \Comment{optional sheaf-based}
\State $\hat{y}_{t+\Delta}, a_t \leftarrow \mathrm{ST\text{-}HyperForecast}(\mathcal{H}_t, c_t)$
\State $\kappa_t \leftarrow \mathrm{QuantumKernelEval}(x_t)$ \Comment{optional regime similarity}
\State $z_t \leftarrow \mathrm{ConstrainedPlan}(\hat{y}_{t+\Delta}, a_t, \kappa_t)$
\State \Return $\hat{y}_{t+\Delta}, a_t, z_t$
\end{algorithmic}
\end{algorithm}

\begin{algorithm}[t]
\caption{Counterfactual / Digital-Twin Scenario Search (Classical + Optional Quantum)}
\begin{algorithmic}[1]
\Require Initial state $\mathcal{H}_t$, intervention family $\mathcal{Z}$, constraints $\mathcal{C}$
\Ensure Counterfactual set $\{(z^{(i)}, \hat{y}^{(i)})\}_{i=1}^N$
\State Initialize proposal distribution over interventions $\pi_0(z)$
\For{$k=0$ to $K-1$}
  \State Sample candidates $z^{(i)} \sim \pi_k$
  \State Score via simulator/forecast $\hat{y}^{(i)} \leftarrow \mathrm{Sim}(\mathcal{H}_t, z^{(i)})$
  \State Filter infeasible candidates violating $\mathcal{C}$
  \State Update $\pi_{k+1}$ toward elite samples (Cross-Entropy Method)
\EndFor
\State \Return best $N$ feasible counterfactuals
\end{algorithmic}
\end{algorithm}

\section{Evaluation Plan}

\subsection{Datasets}
We evaluate on publicly used spatiotemporal forecasting benchmarks (e.g., traffic sensor networks) and
on domain-specific datasets available to partners. Baselines include graph forecasting methods and
hypergraph ablations. \citep{li2018dcrnn,yu2018stgcn,feng2019hypergraph}

\subsection{Metrics}
\begin{itemize}[leftmargin=*]
\item Forecasting: MAE/RMSE/MAPE; calibration and coverage for probabilistic forecasts.
\item Detection: AUROC/AUPRC; time-to-detect for incidents.
\item Planning: constraint violation rate; utility under constraints; robustness under missing data.
\item Compute: latency, energy, memory footprint (edge and server).
\item Governance: DP budget accounting; group disparity metrics; auditability.
\end{itemize}

\begin{table}[t]
\centering
\caption{Benchmark Template (Report with Confidence Intervals)}
\begin{tabular}{lcccc}
\toprule
Metric & Classical ST-GNN & Hypergraph (classical) & QAHUS (no-quantum) & QAHUS (quantum-assisted)\\
\midrule
Traffic MAE (lower better) & TBD & TBD & TBD & TBD \\
Anomaly AUROC (higher better) & TBD & TBD & TBD & TBD \\
Constraint violations (\%) & TBD & TBD & TBD & TBD \\
Edge latency (ms) & TBD & TBD & TBD & TBD \\
Energy per inference (mJ) & TBD & TBD & TBD & TBD \\
\bottomrule
\end{tabular}
\end{table}

\section{Ethical and Regulatory Considerations}
We implement governance by design: minimal raw data retention, differential privacy for releases,
fairness audits, and explicit constrained decision-making. We document model cards, data provenance,
and intervention logs suitable for oversight.

\section{Future Directions}
\begin{itemize}[leftmargin=*]
\item Real-time integration with smart-city APIs and emergency management dashboards.
\item Larger-scale digital twins (multi-million agent simulations) with validated calibration.
\item Hardware-in-the-loop evaluation of quantum subroutines and classical fallbacks.
\item Peer-reviewed dissemination and standardization contributions.
\end{itemize}

\section{Conclusion}
QAHUS turns edge-efficient vision into a city-scale forecasting and constrained decision-support stack
via dynamic hypergraphs. The quantum components target plausible near-term acceleration points
(similarity kernels and discrete optimization) without being single points of failure: the system remains
useful in fully classical mode. By explicitly encoding safety, privacy, and fairness as constraints,
QAHUS aims to be both high-performance and governance-ready.

\bibliographystyle{plainnat}
\bibliography{references}

\end{document}
